% !TeX program = xelatex
% UTF-8; compile twice with XeLaTeX.
% Source: 问题一_定位区域模型.md
% This file preserves the supplied text, mathematics, data and section numbering.
\documentclass[UTF8,a4paper,zihao=-4,fontset=fandol]{ctexart}
\XeTeXgenerateactualtext=1
\usepackage[left=23mm,right=23mm,top=22mm,bottom=23mm,headheight=15pt,headsep=7mm]{geometry}
\usepackage{amsmath,amssymb,mathtools}
\usepackage{booktabs,tabularx,array}
\usepackage{enumitem}
\usepackage{needspace}
\usepackage{tcolorbox}
\usepackage{fancyhdr}
\usepackage[unicode,hidelinks]{hyperref}
\usepackage{bookmark}
\hypersetup{pdftitle={问题一：干扰源定位区域及其直径的确定}}
\renewcommand{\thesection}{5.\arabic{section}}
\renewcommand{\thesubsection}{\thesection.\arabic{subsection}}
\ctexset{
  section={format=\large\bfseries,beforeskip=1.2em,afterskip=0.55em},
  subsection={format=\normalsize\bfseries,beforeskip=0.9em,afterskip=0.35em}
}
\linespread{1.16}
\setlength{\parskip}{0.25em}
\setlength{\emergencystretch}{2em}
\setlist[enumerate]{label=\arabic*.,leftmargin=2em,itemsep=0.15em,topsep=0.4em,parsep=0pt}
\setlist[itemize]{leftmargin=2em,itemsep=0.15em,topsep=0.4em,parsep=0pt}
\pagestyle{fancy}
\fancyhf{}
\fancyhead[L]{\small 问题一：干扰源定位区域}
\fancyfoot[C]{\small\thepage}
\renewcommand{\headrulewidth}{0.3pt}
\renewcommand{\footrulewidth}{0pt}
\fancypagestyle{plain}{\fancyhf{}\fancyfoot[C]{\small\thepage}\renewcommand{\headrulewidth}{0pt}}
\widowpenalty=10000
\clubpenalty=10000
\displaywidowpenalty=10000
\raggedbottom
\begin{document}
\thispagestyle{plain}
\begin{center}{\LARGE\bfseries 问题一：干扰源定位区域及其直径的确定\par}\end{center}
\vspace{0.4em}

\section*{2.1 问题一的分析}
\addcontentsline{toc}{section}{2.1 问题一的分析}

问题一要求根据若干检测点的位置与示向度确定干扰源的可能区域，计算区域的直径，并判断以该直径为直径的圆能否覆盖整个区域。

由于测向存在 $\pm 1^\circ$ 的误差，一次测向只能把干扰源限制在一个张角 $2^\circ$ 的角域内，多次测向的定位区域即各角域的交集。角域可以写成两个线性不等式，于是定位区域是有限个\textbf{半平面的交}，为凸多边形，其直径在顶点对上取得，枚举顶点即可求出。至于覆盖性，我们发现直径为 $D$ 的圆若要盖住区域，圆心只能取最远点对的中点，据此构造出一个符合题设测向条件的算例，其中该圆恰好盖不住区域的一个顶点，从而对第二小问给出否定回答。

\section*{三、模型假设}
\addcontentsline{toc}{section}{三、模型假设}

\begin{enumerate}
\item 测向误差有界，界为 $\pm 1^\circ$，不对误差的分布形式作附加假设；
\item 检测点位置精确已知，问题在平面内讨论；
\item 各测向数据相容，即真实干扰源同时位于所有误差角域内。
\end{enumerate}

\Needspace{21\baselineskip}
\section*{四、符号说明}
\addcontentsline{toc}{section}{四、符号说明}

\begin{center}
\renewcommand{\arraystretch}{1.3}
\begin{tabularx}{0.98\linewidth}{>{\centering\arraybackslash}p{0.25\linewidth}>{\centering\arraybackslash}p{0.24\linewidth}>{\centering\arraybackslash}X}
\toprule
\textbf{符号} & \textbf{范围} & \textbf{说明} \\
\midrule
$n$ & $\mathbb{Z}^+$ & 测向次数 \\
$S_i=(x_i,y_i)$ & $\mathbb{R}^2$ & 第 $i$ 个检测点的坐标 \\
$\theta_i$ & $[0^\circ,360^\circ)$ & 第 $i$ 次测向的示向度 \\
$\varepsilon$ & $1^\circ$ & 测向误差界 \\
$P$ & --- & 定位区域 \\
$v_1,\dots,v_h$ & $\mathbb{R}^2$ & 定位区域的顶点 \\
$D$ & $(0,+\infty)$ & 定位区域的直径 \\
$C^*,\ R^*$ & --- & 最小包围圆的圆心与半径 \\
\bottomrule
\end{tabularx}
\end{center}

\section*{5.1 问题一模型的建立与求解}
\addcontentsline{toc}{section}{5.1 问题一模型的建立与求解}

\subsection*{5.1.1 误差角域与半平面交模型}
\addcontentsline{toc}{subsection}{5.1.1 误差角域与半平面交模型}

示向度自 $x$ 轴正方向逆时针计量，误差不超过 $\pm 1^\circ$，故第 $i$ 次测向的含义是：干扰源位于以 $S_i$ 为顶点、以 $\theta_i$ 为中轴线、两侧各张开 $1^\circ$ 的角域内。记 $\varphi_i=\pi\theta_i/180$，角域两条边界射线的方向为

\[
u_i^-=\begin{pmatrix}\cos(\varphi_i-\varepsilon)\\ \sin(\varphi_i-\varepsilon)\end{pmatrix},\qquad
u_i^+=\begin{pmatrix}\cos(\varphi_i+\varepsilon)\\ \sin(\varphi_i+\varepsilon)\end{pmatrix}.
\]

利用二维叉积 $a\times b=a_xb_y-a_yb_x$ 的符号判断转向，点 $X=(x,y)$ 落在该角域内等价于 $X-S_i$ 夹在两条射线之间：

\[
u_i^-\times(X-S_i)\ge 0,\qquad u_i^+\times(X-S_i)\le 0. \tag{1}
\]

角域张角仅 $2^\circ$，式 (1) 已排除检测点背后的反向区域。展开叉积，每次测向化为两个线性不等式：

\[
\begin{aligned}
\sin(\varphi_i-\varepsilon)(x-x_i)-\cos(\varphi_i-\varepsilon)(y-y_i)&\le 0,\\
-\sin(\varphi_i+\varepsilon)(x-x_i)+\cos(\varphi_i+\varepsilon)(y-y_i)&\le 0.
\end{aligned}
\tag{2}
\]

将全部 $2n$ 个约束统一记为 $a_jx+b_jy\le c_j$，定位区域即

\[
P=\bigcap_{j=1}^{2n}\left\{(x,y):a_jx+b_jy\le c_j\right\}. \tag{3}
\]

半平面之交是凸集，故 $P$ 非空且有界时为凸多边形，退化情形为线段或点。题目图 2 的红色四边形就是 $n=2$ 时两个角域的交。

\subsection*{5.1.2 定位区域及其直径的求解}
\addcontentsline{toc}{subsection}{5.1.2 定位区域及其直径的求解}

求解分三步：判别有界性、提取顶点、枚举最远顶点对。

若各示向度方向接近平行，角域之交会沿测向方向无限延伸，此时谈不上直径。因此先用线性规划检验约束组 (3)：不可行说明测向数据与误差界矛盾；可行时再分别求 $x$、$-x$、$y$、$-y$ 的最大值，四者都有限才说明 $P$ 有界，否则应补充测向后再定位。

$P$ 有界时，其顶点必为某两条边界直线的交点。对直线 $a_jx+b_jy=c_j$ 与 $a_kx+b_ky=c_k$，记 $\Delta_{jk}=a_jb_k-a_kb_j$，当 $\Delta_{jk}\ne 0$ 时交点为

\[
x_{jk}=\frac{c_jb_k-b_jc_k}{\Delta_{jk}},\qquad
y_{jk}=\frac{a_jc_k-c_ja_k}{\Delta_{jk}}. \tag{4}
\]

把交点代回全部约束，保留满足 $a_\ell x_{jk}+b_\ell y_{jk}\le c_\ell$（$\ell=1,\dots,2n$）者并去重，即得顶点 $v_1,\dots,v_h$。

区域直径定义为 $D=\max_{X,Y\in P}\|X-Y\|_2$。对凸多边形，它就等于最远顶点对的距离：

\[
D=\max_{1\le i<j\le h}\|v_i-v_j\|_2 . \tag{5}
\]

这是因为任意 $X,Y\in P$ 都可写成顶点的凸组合 $X=\sum_i\lambda_iv_i$、$Y=\sum_j\mu_jv_j$，于是 $X-Y=\sum_{i,j}\lambda_i\mu_j(v_i-v_j)$ 且系数和为 1，由三角不等式 $\|X-Y\|\le\max_{i,j}\|v_i-v_j\|$；而顶点本身属于 $P$，该上界可以取到。

整个流程归纳如下：

\begin{tcolorbox}[colback=black!2,colframe=black!22,boxrule=0.4pt,arc=0pt,left=9pt,right=9pt,top=7pt,bottom=7pt]
\small\setlength{\parskip}{2pt}
\noindent\hangindent=1.5em\hangafter=1 输入：检测点 $S_i$、示向度 $\theta_i$，$i=1,\dots,n$\par
\noindent\hangindent=1.5em\hangafter=1 1  按式 (2) 将每次测向化为两个半平面；\par
\noindent\hangindent=1.5em\hangafter=1 2  线性规划检验可行性与有界性，异常情形分别报告；\par
\noindent\hangindent=1.5em\hangafter=1 3  按式 (4) 求每对边界直线的交点；\par
\noindent\hangindent=1.5em\hangafter=1 4  保留满足全部约束的交点并去重，得顶点集；\par
\noindent\hangindent=1.5em\hangafter=1 5  枚举所有顶点对，取距离最大者 $(A,B)$，$D=|AB|$。\par
\end{tcolorbox}

边界直线共 $2n$ 条，交点枚举 $O(n^2)$ 次，每个交点检验约束需 $O(n)$，顶点数 $h\le 2n$，故总计算量 $O(n^3)$，对本题的检测点规模完全够用。程序实现时对平行判定、约束检验和去重设置数值容差即可。

\subsection*{5.1.3 直径圆的覆盖性分析}
\addcontentsline{toc}{subsection}{5.1.3 直径圆的覆盖性分析}

设最远顶点对为 $A$、$B$，$\|A-B\|=D$。若某个半径 $D/2$ 的圆盖住了 $P$，它必须同时包含 $A$、$B$，而

\[
D=\|A-B\|\le\|A-C\|+\|C-B\|\le\frac D2+\frac D2=D,
\]

两端相等迫使中间全部取等，故圆心 $C$ 只能是线段 $AB$ 的中点，记为 $C_0$。也就是说，直径恰为 $D$ 的覆盖圆没有挑选圆心的余地；又因为圆盖住全部顶点就盖住整个凸多边形，覆盖性等价于

\[
\max_{1\le i\le h}\|v_i-C_0\|\le\frac D2 . \tag{6}
\]

问题于是变成：式 (6) 是否总成立？我们构造如下算例说明它可以不成立。取两个检测点

\[
S_1=(-500,\,0),\quad\theta_1=44.5^\circ;\qquad
S_2=(500,\,0),\quad\theta_2=135.5^\circ,
\]

误差界仍为 $\pm 1^\circ$。记 $t=\tan 43.5^\circ$、$u=\tan 45.5^\circ$，两个角域分别为 $t(x+500)\le y\le u(x+500)$ 与 $t(500-x)\le y\le u(500-x)$，按 5.1.2 节方法求交得一四边形，结果见表 1。

\Needspace{19\baselineskip}
\begin{center}
\textbf{表 1 算例的顶点与覆盖性检验（单位：m）}
\end{center}

\begin{center}
\renewcommand{\arraystretch}{1.3}
\begin{tabularx}{0.98\linewidth}{>{\raggedright\arraybackslash}p{0.52\linewidth}X}
\toprule
\textbf{量} & \textbf{数值} \\
\midrule
下顶点 $V_1$ & $(0,\ 474.482283)$ \\
右顶点 $V_2$ & $(17.452406,\ 491.043999)$ \\
上顶点 $V_3$ & $(0,\ 508.803696)$ \\
左顶点 $V_4$ & $(-17.452406,\ 491.043999)$ \\
直径 $D=\|V_2-V_4\|$ & $34.904813$ \\
圆心 $C_0=(V_2+V_4)/2$ & $(0,\ 491.043999)$ \\
直径圆半径 $D/2$ & $17.452406$ \\
上顶点到圆心距离 $\|V_3-C_0\|$ & $\boldsymbol{17.759698}$ \\
\bottomrule
\end{tabularx}
\end{center}

最远点对为左右顶点 $V_2$、$V_4$，$D\approx 34.905$ m，圆心只能取 $C_0\approx(0,\ 491.044)$。但上顶点到 $C_0$ 的距离为 $508.803696-491.043999=17.759698$ m，超出半径约 $0.307$ m，落在圆外。

从几何上看原因很清楚：这个四边形两条对角线几乎等长（$34.905$ m 与 $34.321$ m），但水平对角线并不平分竖直对角线---交点之上高 $17.760$ m，之下只有 $16.562$ m。圆心被钉在水平对角线中点，偏低，上顶点自然探出圆外。

该算例由两个角域直接求交产生，符合本题的测向机制；四边形整体位于目标 $1800$ m 圆域内（顶点到原点最远约 $508.8$ m），加入这一先验也不改变结果。故第二小问的答案是否定的：\textbf{以定位区域直径为直径的圆不一定能覆盖该区域，即使允许自由选择圆心也一样}。

\subsection*{5.1.4 保证覆盖时的修正：最小包围圆}
\addcontentsline{toc}{subsection}{5.1.4 保证覆盖时的修正：最小包围圆}

上例说明区域直径 $D$ 与“盖住区域所需的最小圆直径”是两回事。后者由最小包围圆模型给出：

\[
\min_{C,\ R}\ R
\qquad\text{s.t.}\quad\|v_i-C\|_2\le R,\quad i=1,\dots,h, \tag{7}
\]

其中 $C$、$R$ 为待求的圆心与半径；由圆的凸性，约束住顶点即约束住整个区域。有限点集的最小包围圆或以两个顶点为直径两端，或由三个不共线顶点外接确定，顶点不多时枚举这两类候选圆、取覆盖全部顶点且半径最小者即可。对表 1 的算例解得

\[
C^*\approx(0,\ 491.348632),\qquad
R^*\approx 17.455065\ \text{m}>\frac D2\approx 17.452406\ \text{m},
\]

即使圆心取到最优位置，所需半径仍大于 $D/2$。

一般地，平面有界凸区域满足 $\dfrac D2\le R^*\le\dfrac{D}{\sqrt 3}$：下界即最远点对的要求；若最小包围圆由三点外接确定，这三点构成锐角三角形，最大内角 $\alpha\in[60^\circ,90^\circ)$，由正弦定理 $R^*=\dfrac{\text{对边}}{2\sin\alpha}\le\dfrac{D}{\sqrt 3}$，等边三角形取到上界。可见把覆盖圆半径取成 $D/2$ 最多会低估约 $15.5\%$。

这一区分与后续问题直接相关：无人机须进入距干扰源 $20$ m 范围才能光学定位，要保证在某一位置能覆盖全部候选位置，应检验 $R^*\le 20$ m，而 $D\le 40$ m 只是必要条件。

\subsection*{5.1.5 问题一小结}
\addcontentsline{toc}{subsection}{5.1.5 问题一小结}

$n$ 次测向的定位区域是 $2n$ 个半平面的交，非空有界时为凸多边形；枚举边界直线交点得到顶点后，区域直径即最远顶点对的距离，总计算量 $O(n^3)$。以 $D$ 为直径的覆盖圆圆心只能取最远点对中点，而构造的算例（$D\approx 34.905$ m）中该圆漏掉上顶点约 $0.307$ m，最小包围圆半径 $R^*\approx 17.455$ m 也大于 $D/2$，故直径圆不一定能覆盖定位区域；需要保证覆盖时应改用最小包围圆，并以 $R^*\le 20$ m 作为光学定位可达的判据。
\end{document}
